\documentclass[11pt]{article}
\usepackage[margin=1in]{geometry}
\usepackage{amsmath,amssymb,amsthm,mathtools}
\usepackage{graphicx}
\usepackage[colorlinks=true,linkcolor=blue,citecolor=blue,urlcolor=blue]{hyperref}

\newtheorem{theorem}{Theorem}
\newtheorem{lemma}{Lemma}
\newtheorem{proposition}{Proposition}
\newtheorem{corollary}{Corollary}
\newtheorem{definition}{Definition}
\newtheorem{remark}{Remark}

\newcommand{\E}{\mathsf E}
\newcommand{\Ed}{\mathsf E_\delta}
\newcommand{\Szero}{\mathsf S_0}
\newcommand{\supp}{\operatorname{supp}}
\newcommand{\Bor}{\operatorname{Bor}}
\newcommand{\N}{\mathbb N}
\newcommand{\R}{\mathbb R}
\newcommand{\eps}{\varepsilon}

\title{Zero-Dimensional Measure Extraction in Compact Spaces}
\author{Run \texttt{fa\_banach\_001}: consolidated stronger partial packet}
\date{18 June 2026}

\begin{document}
\maketitle

\section*{Source Problem and Status}

The source paper is Damian Glodkowski, \emph{A Banach space \(C(K)\) reading
the dimension of \(K\)}, arXiv:2207.00149, later published in
\emph{Journal of Functional Analysis}.  Problem 7.3 asks:

\begin{quote}
\textbf{Problem 7.3.} Describe the class of compact Hausdorff spaces \(K\) such
that for every non-zero Radon measure \(\mu\) on \(K\) there is a
zero-dimensional compact subset \(L\subseteq K\) such that \(\mu(L)\neq 0\).
\end{quote}

\begin{center}
\includegraphics[width=0.94\linewidth]{figures/open_problem_crop.png}
\end{center}

\textbf{Verdict.}  This is a substantially strengthened partial answer, not a
complete classification of all compact Hausdorff spaces.  The previous packet
proved finite-product and restricted-extension closure for a property that it
called ``strong''.  The corrected result below shows that this strong-looking
approximation property is actually equivalent to the original positive
extraction property.  It also proves countable closure, removes the unnecessary
exceptional-set hypothesis from the source-paper extension step, gives ZFC
nonexamples, and proves an exact characterization for arbitrary products.

\section{Definitions}

All compact spaces are Hausdorff, all measures are finite Radon measures, and
``zero-dimensional'' means covering dimension zero.  For compact Hausdorff
spaces this is equivalent to having a base of clopen sets.

\begin{definition}
A compact space \(K\) belongs to \(\E\) if every nonzero finite positive Radon
measure \(\mu\) on \(K\) has a compact zero-dimensional subset \(Z\subseteq K\)
with \(\mu(Z)>0\).
\end{definition}

\begin{definition}
A compact space \(K\) belongs to \(\Ed\) if every nonzero finite positive Radon
measure \(\mu\) on \(K\) has a compact zero-dimensional \(G_\delta\) subset
\(Z\subseteq K\) with \(\mu(Z)>0\).
\end{definition}

The class \(\Ed\) is the version relevant to the countable-neighborhood
extraction used in the metric part of the source paper.  The class \(\E\) is
the literal formulation of Problem 7.3 for positive measures.

\section{Extraction is Equivalent to Approximation}

\begin{theorem}[Extraction, approximation, and decomposition]
For a compact space \(K\), the following are equivalent.
\begin{enumerate}
\item \(K\in\E\).
\item For every finite positive Radon measure \(\mu\) on \(K\) and every
\(\eta>0\), there is a compact zero-dimensional \(Z\subseteq K\) such that
\[
  \mu(K\setminus Z)<\eta .
\]
\item Every finite positive Radon measure \(\mu\) on \(K\) can be written
\[
  \mu=\sum_{n=1}^{\infty}\mu_n
\]
as an increasing total-variation sum of positive Radon measures such that
\(\supp\mu_n\) is zero-dimensional for every \(n\).
\end{enumerate}
\end{theorem}

\begin{proof}
Assume \(K\in\E\), and fix a finite positive measure \(\mu\).  Let
\[
 s=\sup\{\mu(Z):Z\subseteq K\text{ compact and zero-dimensional}\}.
\]
Choose compact zero-dimensional \(Z_n\) with \(\mu(Z_n)>s-2^{-n}\), replacing
\(Z_n\) by \(Z_1\cup\cdots\cup Z_n\).  Finite closed sums of zero-dimensional
compacta are zero-dimensional, so \(Z_n\) remains compact zero-dimensional and
the sequence is increasing.  If \(s<\mu(K)\), let \(B=K\setminus\bigcup_n Z_n\)
and define the positive Radon measure \(\lambda(A)=\mu(A\cap B)\) on Borel
sets.  By \(\E\), there is compact zero-dimensional \(W\) with
\(\lambda(W)>0\).  For all sufficiently large \(n\),
\[
  \mu(Z_n\cup W)\ge \mu(Z_n)+\lambda(W)>s,
\]
contradicting the definition of \(s\).  Thus \(s=\mu(K)\), which proves the
approximation statement.

Assume the approximation statement.  Choose compact zero-dimensional
\(W_n\subseteq K\) with \(\mu(K\setminus W_n)<2^{-n}\), and put
\(Z_n=W_1\cup\cdots\cup W_n\).  With \(Z_0=\emptyset\), \(B_n=Z_n\setminus
Z_{n-1}\), and \(\mu_n(A)=\mu(A\cap B_n)\), each \(\mu_n\) is a positive Radon
measure and \(\supp\mu_n\subseteq Z_n\).  Since \(\mu(K\setminus\bigcup_n
Z_n)=0\), we have \(\mu=\sum_n\mu_n\) in total variation.

Finally, if \(\mu=\sum_n\mu_n\neq0\), then some \(\mu_n\neq0\).  Its support is
compact zero-dimensional and has positive \(\mu_n\)-measure, hence positive
\(\mu\)-measure.  This gives \(\E\).
\end{proof}

\begin{remark}
The same supremum argument proves the analogous equivalence for \(\Ed\): it is
enough to replace ``compact zero-dimensional'' by ``compact zero-dimensional
\(G_\delta\)''.  Finite unions of \(G_\delta\) sets are \(G_\delta\).
\end{remark}

\begin{proposition}[Signed and complex measures]
The positive-measure formulation of \(\E\) is equivalent to the original
Problem 7.3 formulation for real signed Radon measures, and also to the
analogous formulation for complex Radon measures.
\end{proposition}

\begin{proof}
Let \(\nu\) be a nonzero signed measure.  Replacing \(\nu\) by \(-\nu\) if
necessary, \(\nu^+\neq0\).  Choose a Hahn-positive Borel set \(P\), then by
inner regularity a compact \(F\subseteq P\) with \(\nu^+(F)>0\).  Applying
\(\E\) inside \(F\) to \(\nu^+|F\) gives compact zero-dimensional
\(Z\subseteq F\) with \(\nu^+(Z)>0\), and \(\nu^-(Z)=0\).  Hence
\(\nu(Z)>0\).  For a complex measure \(\nu\), choose \(\theta\) so that
\(\operatorname{Re}(e^{-i\theta}\nu)\) is nonzero and apply the real case.
Here we use the elementary heredity of \(\E\) for compact subspaces, proved
explicitly in the closure theorem below.
\end{proof}

\section{Metric Spaces, Closure, and the Source Extensions}

\begin{theorem}[Compact metrizable spaces]
Every compact metrizable space belongs to \(\Ed\).
\end{theorem}

\begin{proof}
Let \(K\) be compact metrizable, let \(\mu\) be a finite positive measure, and
let \(\eta>0\).  Fix a compatible metric, a countable dense set
\((d_n)\), and \(D=\operatorname{diam}K\).  For each \(n\), the spheres
\[
  S(d_n,r)=\{x:d(x,d_n)=r\},\qquad 0<r<D,
\]
are pairwise disjoint, so only countably many have positive \(\mu\)-measure.
Choose distinct radii \(r_{n,j}\) outside that exceptional set so that every
tail \(\{r_{n,j}:j\ge k\}\) is dense in \((0,D)\).  Put
\[
  Y=K\setminus\bigcup_{n,j} S(d_n,r_{n,j}).
\]
Then \(\mu(K\setminus Y)=0\).  The sets \(B(d_n,r_{n,j})\cap Y\) are clopen in
\(Y\) and form a base: given \(y\in O\), choose \(\rho>0\) with
\(B(y,\rho)\subseteq O\), choose \(d_n\) within \(\rho/3\) of \(y\), and then
choose a tail radius in \((d(y,d_n),\rho-d(y,d_n))\).  Thus \(Y\) is
zero-dimensional.  By inner regularity choose compact \(Z\subseteq Y\) with
\(\mu(K\setminus Z)<\eta\).  Since \(K\) is metrizable, \(Z\) is a \(G_\delta\)
subset of \(K\).  Therefore \(K\in\Ed\).
\end{proof}

\begin{theorem}[Countable closure]
The following closure properties hold.
\begin{enumerate}
\item \(\E\) is hereditary for compact subspaces, closed under countable closed
covers, and closed under countable products.
\item \(\Ed\) is hereditary for compact subspaces and closed under countable
products.
\item Both \(\E\) and \(\Ed\) are closed under inverse limits indexed by a
countable set, in particular by a countable ordinal.
\end{enumerate}
\end{theorem}

\begin{proof}
For heredity, let \(F\subseteq K\) be compact and extend a measure on \(F\) to
a measure on \(K\) supported by \(F\).  Intersecting the extracting compactum
with \(F\) gives the result; in the \(G_\delta\) case the intersection is
\(G_\delta\) in \(F\).

If \(K=\bigcup_n F_n\), where the \(F_n\) are compact and belong to \(\E\), then
any nonzero positive measure on \(K\) has \(\mu(F_n)>0\) for some \(n\), so
extraction inside \(F_n\) proves \(K\in\E\).

Let \(K=\prod_{n\in\N}K_n\), with each \(K_n\in\E\), and let \(\mu\) be a
positive probability measure on \(K\).  For the marginal
\(\mu_n=(p_n)_*\mu\), choose compact zero-dimensional \(A_n\subseteq K_n\) with
\[
  \mu_n(K_n\setminus A_n)<\eta 2^{-n-1}.
\]
Then \(A=\prod_n A_n\) is compact zero-dimensional and
\[
  \mu(K\setminus A)\le\sum_n\mu_n(K_n\setminus A_n)<\eta .
\]
This gives \(\E\).  If each \(K_n\in\Ed\), choose each \(A_n\) to be \(G_\delta\);
then \(A=\bigcap_n p_n^{-1}(A_n)\) is \(G_\delta\) in \(K\).

Finally, a countable inverse limit of compact Hausdorff spaces is a closed
subspace of their countable product.  Apply product closure and heredity.
\end{proof}

\begin{corollary}[Source-paper extensions]
Let \(L\) be any extension of \(K\) by a pairwise disjoint sequence of
\([0,1]\)-valued continuous functions in the sense of the source paper.  Then
\[
  K\in\E \Rightarrow L\in\E,\qquad K\in\Ed \Rightarrow L\in\Ed .
\]
No strongness assumption and no hypothesis on \(K\setminus D((f_n))\) is needed.
\end{corollary}

\begin{proof}
By the definition used in the source paper, \(L\) is a compact subspace, hence a
closed subspace, of \(K\times[0,1]\).  The interval belongs to \(\Ed\) by the
metric theorem.  Product closure and heredity give the result.
\end{proof}

\begin{corollary}
Starting from a compactum in \(\E\), respectively in \(\Ed\), every countable
tower formed by source-paper extensions at successor stages and inverse limits
at countable limit stages remains in \(\E\), respectively in \(\Ed\).
\end{corollary}

\begin{remark}
This repairs the earlier packet's \(G_\delta\)/countable-character concern for
countable stages.  It still does not remove all metric input from the source
paper's construction: the proof of Proposition 6.2 uses a weak-star convergent
subsequence of a bounded sequence of measures, and later selection arguments
also use metric structure.
\end{remark}

\section{ZFC Nonexamples}

\begin{proposition}[Normal-measure counterexample]
If a nondegenerate compact connected space \(K\) carries a normal probability
measure, meaning a measure that vanishes on every Borel set with empty interior,
then \(K\notin\E\).
\end{proposition}

\begin{proof}
Let \(F\subseteq K\) be compact zero-dimensional.  If \(F\) had nonempty
interior, choose \(x\in\operatorname{int}F\) and \(y\in K\setminus\{x\}\).  A
clopen base in \(F\) gives a nonempty set \(C\) clopen in \(F\) with
\[
 x\in C\subseteq(\operatorname{int}F)\setminus\{y\}.
\]
Since \(F\) is closed and \(C\subseteq\operatorname{int}F\), the set \(C\) is
both closed and open in \(K\), and it is proper.  This contradicts connectedness.
Thus every compact zero-dimensional \(F\) has empty interior, so a normal
probability measure assigns measure zero to all such \(F\).
\end{proof}

Plebanek constructed in ZFC a nondegenerate compact connected space carrying a
normal probability measure.  Hence not every compact Hausdorff space belongs to
\(\E\).

\section{Exact Characterization for Arbitrary Products}

\begin{definition}
A compact space \(K\) belongs to \(\Szero\) if \(\supp\mu\) is zero-dimensional
for every finite positive Radon measure \(\mu\) on \(K\).
\end{definition}

\begin{lemma}[Supports]
For a positive Radon measure \(\mu\) on compact \(K\), \(\supp\mu\) is compact,
has full \(\mu\)-measure, and every nonempty relatively open subset of
\(\supp\mu\) has positive \(\mu\)-measure.
\end{lemma}

\begin{proof}
Let \(U\) be the union of all open \(\mu\)-null subsets of \(K\).  Every compact
subset of \(U\) is covered by finitely many open null sets, hence is null; inner
regularity gives \(\mu(U)=0\).  Thus \(K\setminus U=\supp\mu\) is compact and
has full measure.  The final assertion is immediate from the definition of
support.
\end{proof}

\begin{lemma}
The class \(\Szero\) is hereditary for compact subspaces and closed under
arbitrary products.
\end{lemma}

\begin{proof}
Heredity follows by extending a measure from a compact subspace to the ambient
space; the subspace support is a closed subspace of the ambient support.

Let \(K=\prod_{i\in I}K_i\), with each \(K_i\in\Szero\), and let \(\mu\) be a
positive measure on \(K\).  Put \(\mu_i=(p_i)_*\mu\) and \(S_i=\supp\mu_i\).
If \(x\in\supp\mu\), then \(x_i\in S_i\) for every \(i\), since an open
\(\mu_i\)-null neighborhood of \(x_i\) would pull back to an open \(\mu\)-null
neighborhood of \(x\).  Hence
\[
  \supp\mu\subseteq\prod_{i\in I}S_i .
\]
Each \(S_i\) is zero-dimensional.  Arbitrary products of compact
zero-dimensional spaces are zero-dimensional, and closed subspaces preserve
zero-dimensionality.  Thus \(\supp\mu\) is zero-dimensional.
\end{proof}

\begin{lemma}[Positive compact subsets contain full tail fibers]
Let \(X=\prod_{i\in I}X_i\), where each \(X_i\) is nonempty compact and carries a
strictly positive Radon probability \(\lambda_i\).  Let
\(\lambda=\bigotimes_{i\in I}\lambda_i\) be the canonical product probability.
If \(F\subseteq X\) is compact and \(\lambda(F)>0\), then there are a countable
set \(C\subseteq I\) and a point \(a\in\prod_{i\in C}X_i\) such that
\[
  \{a\}\times\prod_{i\in I\setminus C}X_i\subseteq F .
\]
\end{lemma}

\begin{proof}
Let \(U=X\setminus F\).  For each \(n\), choose compact \(H_n\subseteq U\) with
\(\lambda(H_n)>\lambda(U)-2^{-n}\).  Cover \(H_n\) by finitely many basic open
cylinders contained in \(U\).  Let \(V\) be the union of all cylinders obtained
in this way.  Then \(V\subseteq U\), \(V\) depends on only countably many
coordinates, and \(\lambda(V)=\lambda(U)\).

Thus \(V=p_C^{-1}(B)\) for some countable \(C\subseteq I\) and open
\(B\subseteq X_C=\prod_{i\in C}X_i\).  Put \(A=X_C\setminus B\) and
\(P=p_C^{-1}(A)\).  Then \(F\subseteq P\) and \(\lambda(P\setminus F)=0\).
Writing \(T=I\setminus C\), let \(\lambda_C\) and \(\lambda_T\) be the
corresponding product measures and let \(\nu=\lambda_C|A\).  Since
\(\lambda(P)\ge\lambda(F)>0\), \(\nu\neq0\).  Choose \(a\in\supp\nu\).

The measure \(\lambda_T\) is strictly positive, so
\[
  \supp(\nu\otimes\lambda_T)=\supp\nu\times X_T .
\]
The relatively open set \(P\setminus F\) is null for \(\nu\otimes\lambda_T\),
hence it misses this support.  Therefore \(\{a\}\times X_T\subseteq F\).
\end{proof}

\begin{theorem}[Arbitrary product characterization]
Let \((K_i)_{i\in I}\) be nonempty compact spaces and put
\[
  K=\prod_{i\in I}K_i .
\]
Then \(K\in\E\) if and only if both conditions hold:
\begin{enumerate}
\item \(K_i\in\E\) for every \(i\in I\);
\item the set \(J=\{i\in I:K_i\notin\Szero\}\) is countable.
\end{enumerate}
\end{theorem}

\begin{proof}
Assume the two conditions and let \(\mu\) be a nonzero positive measure on \(K\).
For each \(i\), put \(\mu_i=(p_i)_*\mu\) and \(S_i=\supp\mu_i\).  By the support
lemma, \(\supp\mu\) has full measure and
\[
  \supp\mu\subseteq P:=\prod_{i\in I}S_i .
\]
If \(i\notin J\), then \(S_i\) is zero-dimensional.  If \(i\in J\), then
\(S_i\in\E\) by heredity from \(K_i\).  Since \(J\) is countable,
\(\prod_{i\in J}S_i\in\E\).  The product over \(I\setminus J\) is
zero-dimensional, hence belongs to \(\E\); a finite product with a
zero-dimensional compactum preserves \(\E\).  Therefore \(P\in\E\).  Applying
\(\E\) to the restriction of \(\mu\) to the full-measure compact set \(P\) gives
the required positive-measure zero-dimensional compact subset of \(K\).

Conversely, assume \(K\in\E\).  Each factor \(K_i\) is homeomorphic to a compact
subspace of \(K\) after fixing base points in the other coordinates, so
\(K_i\in\E\).  If \(J\) were uncountable, choose for each \(i\in J\) a
probability \(\mu_i\) on \(K_i\) whose support \(S_i\) is not zero-dimensional.
Regard \(\mu_i\) as a strictly positive probability on \(S_i\), and form
\[
  X=\prod_{i\in J}S_i,\qquad \lambda=\bigotimes_{i\in J}\mu_i .
\]
Let \(F\subseteq X\) be compact with \(\lambda(F)>0\).  The tail-fiber lemma
gives a countable \(C\subseteq J\) and a full tail fiber contained in \(F\).
Choose \(i\in J\setminus C\).  By fixing all tail coordinates except the
\(i\)-th one, \(F\) contains a closed copy of \(S_i\).  Since \(S_i\) is not
zero-dimensional, \(F\) is not zero-dimensional.  Hence \(\lambda\) vanishes on
every compact zero-dimensional subset of \(X\), so \(X\notin\E\).  After fixing
one base point in every coordinate \(i\in I\setminus J\), the product \(X\) is
homeomorphic to a compact subspace of \(K\), contradicting heredity.  Thus \(J\)
is countable.
\end{proof}

\begin{corollary}[Products of compact metrizable spaces]
For nonempty compact metrizable spaces \((M_i)_{i\in I}\),
\[
  \prod_{i\in I}M_i\in\E
  \quad\Longleftrightarrow\quad
  \{i\in I:\dim M_i>0\}\text{ is countable}.
\]
In particular,
\[
  [0,1]^I\in\E \quad\Longleftrightarrow\quad I\text{ is countable}.
\]
\end{corollary}

\begin{proof}
Every compact metrizable space belongs to \(\E\).  Also, a compact metrizable
space \(M\) belongs to \(\Szero\) if and only if it is zero-dimensional: if
\((x_n)\) is dense in \(M\), then \(\sum_n2^{-n}\delta_{x_n}\) is a strictly
positive probability measure with support \(M\).  Apply the product
characterization.
\end{proof}

\begin{corollary}
If \(I\) is uncountable, the canonical product Lebesgue probability on
\([0,1]^I\) assigns measure zero to every compact zero-dimensional subset.
\end{corollary}

\begin{proof}
A compact subset of positive product measure would contain a full tail fiber
outside countably many coordinates.  Since \(I\) is uncountable, that fiber
contains a closed copy of \([0,1]\), so the compact subset is not
zero-dimensional.
\end{proof}

\begin{remark}
The cube \([0,1]^{\omega_1}\) is a sharp obstruction to replacing countable
inverse-limit closure by \(\omega_1\)-length closure: it is the inverse limit of
the compact metrizable cubes \([0,1]^\alpha\) for \(\alpha<\omega_1\), but it
does not belong to \(\E\).
\end{remark}

\section*{Novelty and Human Review Notes}

The supplied audit report states that the exact problem is repeated as Problem
3.5.3 in Glodkowski's 2024 doctoral thesis and that no later complete
classification was found in searches through 18 June 2026.  A quick check by
the current agent also located the source paper but not a later public paper
answering Problem 7.3.  The packet should therefore be reviewed as a likely new
strong partial answer, with the exact product theorem as the most substantive
piece.

The highest-value human checks are:
\begin{enumerate}
\item the proof that positive extraction implies arbitrary good compact
zero-dimensional approximation;
\item the tail-fiber lemma for arbitrary product Radon measures;
\item the necessity direction of the arbitrary product characterization;
\item the claim that source-paper extensions are always closed subspaces of
\(K\times[0,1]\) under Definition 5.1 of the source paper.
\end{enumerate}

\section*{References}

\begin{enumerate}
\item M. Dzamonja and K. Kunen, \emph{Properties of the class of measure
separable compact spaces}, Fund. Math. 147 (1995), 261--277.
\item R. Engelking, \emph{Theory of Dimensions, Finite and Infinite},
Heldermann, 1995.
\item D. Glodkowski, \emph{A Banach space \(C(K)\) reading the dimension of
\(K\)}, J. Funct. Anal. 285 (2023), Paper No. 109986; arXiv:2207.00149.
\item D. Glodkowski, \emph{Some Applications of Set Theory in Banach Spaces and
Operator Algebras}, Ph.D. thesis, Institute of Mathematics of the Polish
Academy of Sciences, Warsaw, April 2024.
\item E. Hewitt, \emph{A remark on density characters}, Bull. Amer. Math. Soc.
52 (1946), 641--643.
\item G. Plebanek, \emph{A normal measure on a compact connected space},
preprint, arXiv:1507.02845.
\item G. Plebanek, \emph{Musing on Kunen's compact L-space}, preprint,
arXiv:2012.01849.
\item O. Zindulka, \emph{Dimension zero vs. measure zero}, Proc. Amer. Math.
Soc. 128 (2000), 1769--1778.
\end{enumerate}

\end{document}
